\subsection{Scope among del Pezzo cones}\label{sec01:degree-scope}

The degree-seven and degree-six cases arise from the pentagonal and
hexagonal faces in the toric examples.  Both associated cones are smoothable and are
not complete intersections.  Their local deformation spaces exhibit
the two phenomena treated here: distinct smoothing components in
degree six, and a nonreduced structure transverse to the smoothing
line in degree seven.  Theorems~\ref{thm:criterionintro}
and~\ref{thm:smoothintro} also allow the degree-five cone and the
anticanonical cone over $\PP^1\times\PP^1$, of degree eight.
These hypotheses specify the
local comparisons established below; smoothable del Pezzo cones also
occur in the other degrees.

For a smooth del Pezzo surface $E$, the anticanonical cone means
\[
 \operatorname{Spec}\bigoplus_{m\geq0}H^0(E,-mK_E).
\]
This convention includes the weighted
anticanonical models in degrees one and two.  The classical local
facts in Table~\ref{tab01:degree-scope} are recalled in
\cite[\S\S3.3, 5.5 and Example~2.19]{Gross97rev}; the branch structures
in degrees six and seven follow from \cite{Altmann97}.
Proposition~\ref{prop03:degree5family} proves the full local-base
and discriminant statements in degree five.

% Canonical degree-scope table.  Keep its final column synchronized with
% the theorem statements and the degree-five/degree-eight research TODO.
% Local smoothability and an established global criterion are distinct.
\begin{table}[htbp]
\centering
\caption{Anticanonical cones over smooth del Pezzo surfaces and the
scope of the present results.  The local assertions alone do not imply
smoothability of a compact threefold containing these germs.}
\label{tab01:degree-scope}
\small
\renewcommand{\arraystretch}{1.14}
\begin{adjustbox}{max width=\textwidth, center}
\begin{tabular}{@{}>{\raggedright\arraybackslash}p{0.17\textwidth}
                  >{\raggedright\arraybackslash}p{0.41\textwidth}
                  >{\raggedright\arraybackslash}p{0.32\textwidth}@{}}
\toprule
Degree / surface & Local deformation behavior & Role in the present results \\
\midrule
$1$--$3$
  & Hypersurfaces; smoothable, with smooth local bases.
  & Classical complete-intersection cases; outside the stated mixed criterion. \\
\addlinespace
$4$
  & Intersection of two quadrics; smoothable, with smooth local base.
  & As in degrees one to three. \\
\addlinespace
$5$
  & Smooth four-dimensional base; five discriminant hyperplanes.
  & Covered by Theorems~\ref{thm:criterionintro}
    and~\ref{thm:smoothintro}; five nonzero conic-pencil pairings. \\
\addlinespace
$6$
  & A smoothing line and a smoothing plane; three discriminant lines in the plane.
  & Both components covered by Theorems~\ref{thm:criterionintro}
    and~\ref{thm:smoothintro}. \\
\addlinespace
$7$
  & One reduced smoothing line, with an obstructed transverse tangent.
  & Covered by Theorems~\ref{thm:criterionintro}
    and~\ref{thm:smoothintro}. \\
\addlinespace
$8$, $\PP^1\!\times\!\PP^1$
  & Smooth one-dimensional local base; noncentral fibres smooth.
  & Covered by Theorems~\ref{thm:criterionintro}
    and~\ref{thm:smoothintro}; coefficient on the ruling difference. \\
\addlinespace
$8$, $\mathbb F_1$
  & No smoothing; local base $\operatorname{Spec}\CC[t]/(t^2)$.
  & Excluded by local nonsmoothability. \\
\addlinespace
$9$, $\PP^2$
  & Rigid.
  & Excluded by local nonsmoothability. \\
\bottomrule
\end{tabular}
\end{adjustbox}
\end{table}

The degree-eight cone over $\PP^1\times\PP^1$ already occurs as the
partial resolution of a degree-seven point.
Lemma~\ref{lem03:quadric} proves smoothness of its full local base and
identifies its tangent line with the difference of the ruling classes.
When such points occur on the original threefold, the partial
modification is the identity near them.  Their Milnor fibres retract
to $\mathbb{RP}^3$; the descended nodal period supplies the necessity
argument in Lemma~\ref{lem04:periodnecessity}.  Thus they require one
nonzero coefficient on the ruling difference, with no additional
hypothesis in the selected-base smoothness theorem.

The degree-five surface is not toric.  Totaro's Bott vanishing
theorem \cite{totaro2020} supplies its local Hodge comparison.
The four-parameter family of Grassmannian sections and its
boundary periods identify the five conic-pencil conditions and
prove necessity at every order.  Degree-five points also remain
unchanged on the partial model.  Thus the present criterion covers
all locally smoothable anticanonical cones over smooth del Pezzo
surfaces that are not complete intersections.
For degrees one to four, the existing complete-intersection theory
does not by itself give the stated criterion for mixtures with the
non-complete-intersection cones considered in this paper.
